A classifier's training dynamics can reveal how it learns from noisy
labels, but new images have no measured training history. We study two
ways to transfer early-training information to inference with frozen
ResNet50 features. The first predicts a class-conditioned probability
margin averaged over five early epochs and uses the predicted margin to
weight test-time augmentations. In exploratory experiments on three
independently trained bases per dataset, this scalar method gains 0.82
percentage points over uniform augmentation on CIFAR-100N and 1.07 on
CIFAR-10N Worse. The corresponding gains are 0.25 on Animal-10N and 0.75
on Mini-WebVision. Calibration does not consistently improve.

The second approach distills a five-checkpoint teacher into a one-pass
posterior head. On a prespecified CIFAR-10.1 test, residual distillation
changes accuracy by $-0.07$ percentage points against the base
(conditional paired 95\% interval $[-0.53, 0.42]$) and meets the latency
criterion. On CIFAR-10.2, the prespecified comparison of direct students
from the five-checkpoint and five-view teachers gives $-0.02$ points
($[-0.53, 0.48]$). In the secondary comparison the five-checkpoint teacher
itself gains 2.10 points over the five-view teacher. The exploratory
scalar results and the external tests characterize different uses of
early-training information. The tested one-pass heads did not meet
either prespecified primary accuracy criterion.
